\documentclass[USenglish, thumbmarks]{uiophdthesis} % https://www.uio.no/english/about/designmanual/profile-in-use/latex/
\usepackage[utf8]{inputenc}

\usepackage[T1]{url}
\urlstyle{same} % no special font (template default is sf)
\usepackage[hidelinks, hypertexnames=false]{hyperref}

\usepackage{babel}
\usepackage{csquotes}
\usepackage{graphicx}
\usepackage{textcomp}
\usepackage{bookmark}

\usepackage[table]{xcolor} % color tables
\usepackage{booktabs}

\usepackage[nospace]{varioref}
\usepackage{cleveref}

\usepackage[backend=biber, style=numeric-comp, url=false]{biblatex}
\DeclareFieldFormat{titlecase}{#1}

\newcommand\TODO[1]{\textcolor{red}{\textbf{TODO:} #1}}
\newcommand\arxiv[1]{\href{https://arxiv.org/abs/#1}{arXiv: #1}}

\title{New tools for exploring modified cosmological models}
%\subtitle{Subtitle}
\author{Herman Sletmoen}

\addbibresource{thesis.bib}

\begin{document}

\frontmatter{}
\maketitle[
  supervisor={Hans A. Winther},
  dept={Institute of Theoretical Astrophysics},
  %fac={Faculty of Mathematics and Natural Sciences},
  %info={Additional information},
  %year={2026}
]

\begin{abstract}
\TODO{3--6 sentences describing your thesis.}
\end{abstract}
\begin{xabstract}[Sammendrag]
\TODO{In Norwegian}
\end{xabstract}

\tableofcontents{}
\listoffigures{}
\listoftables{}

\begin{preface}
\TODO{Preface, including acknowledgments and thanks.}
\section*{Acknowledgements}
\TODO{}
\end{preface}

\mainmatter{}

\chapter{Background}
\section{The state of cosmology today}

\TODO{}

\cite{sletmoenSimplePredictionNonlinear2024}

\cite{sletmoenSymBoltzjlSymbolicnumericApproximationfree2026}

\section{The motivation and development of SymBoltz}

SymBoltz repository \url{https://github.com/hersle/SymBoltz.jl}

\begin{table}[htbp]
\begin{center}
\small
\setlength{\tabcolsep}{2pt}
\def\arraystretch{1.5}
\begin{tabular}{lllllll}
\toprule
 & CAMB & CLASS & PyCosmo & DISCO-EB & Bolt${}^1$ & SymBoltz \\
\midrule
Best implicit solver  & \cellcolor{red}None           & \cellcolor{green}\texttt{ndf15}    & \cellcolor{green}\texttt{BDF2}  & \cellcolor{green}\texttt{Kvaerno5} & \cellcolor{green}\texttt{KenCarp4} & \cellcolor{green}\texttt{Rodas5P} \\
Approximation-free    & \cellcolor{red}No             & \cellcolor{yellow}Almost${}^2$     & \cellcolor{green}Yes            & \cellcolor{green}Yes               & \cellcolor{green}Yes               & \cellcolor{green}Yes \\
Solver order (stable) & 6                             & \cellcolor{yellow}1--5 (2)${}^3$   & \cellcolor{yellow}2 (2)         & \cellcolor{green}5 (5)             & \cellcolor{green}4 (4)             & \cellcolor{green}5 (5) \\
Newton iterations     & --                            & \cellcolor{red}Yes                 & \cellcolor{green}No${}^4$       & \cellcolor{red}Yes                 & \cellcolor{red}Yes                 & \cellcolor{green}No${}^5$ \\
Jacobian method       & --                            & \cellcolor{red}$O(N)$ FD           & \cellcolor{green}$O(1)$ anal.   & \cellcolor{yellow}$O(N)$ AD        & \cellcolor{yellow}$O(N)$ AD        & \cellcolor{green}$O(1)$ anal. \\
Jacobian sparsity     & --                            & \cellcolor{yellow}Changes${}^6$    & \cellcolor{green}Fixed          & \cellcolor{red}None                & \cellcolor{red}None                & \cellcolor{green}Fixed \\
Differentiable        & \cellcolor{red}No             & \cellcolor{red}No                  & \cellcolor{red}No               & \cellcolor{green}Yes               & \cellcolor{green}Yes               & \cellcolor{green}Yes \\
\bottomrule
\end{tabular}
\end{center}
\caption{%
Characteristic of SymBoltz compared to other Einstein--Boltzmann codes.
AD or FD means that the Jacobian is computed with automatic differentiation or finite differences, respectively.
${}^1$Development appears to have ceased.
${}^2$Tight-coupling approximation is mandatory, but others can be disabled (\url{https://github.com/lesgourg/class_public/blob/e85808324f51fc694d12e3ed7439552a3c3f9540/include/perturbations.h\#L42}).
${}^3$BDF methods are \enquote{A-stable} and \enquote{L-stable} only to 2nd order (e.g. \url{https://doi.org/10.1137/S1064827594276424}).
${}^4$Custom `BDF2` method specializing on linearity $\mathbf{f} = \mathbf{J} \mathbf{u}$ (see \url{https://arxiv.org/abs/1708.05177} eq. (21)).
${}^5$Rosenbrock methods linearize $\mathbf{f} = \mathbf{J} \mathbf{u}$ (e.g. \url{https://doi.org/10.1007/s10543-023-00967-x}).
${}^6$Sparsity found numerically and changes with approximations.
}
\label{tab:boltzmann-comparison}
\end{table}

\uiopaper[%compact,
  authors={Herman Sletmoen\and Hans A. Winther},
  published={Published in \href{https://doi.org/10.1051/0004-6361/202450050}{Astronomy \& Astrophysics (Sept. 2024)}},
  info={Equivalent preprint on \arxiv{2403.03786}},
]{A simple prediction of the nonlinear matter power spectrum in Brans–Dicke gravity from linear theory}
\uioincludepdf{paper1.pdf}

\uiopaper[%compact,
  authors={Herman Sletmoen},
  published={Published in \href{https://doi.org/10.1051/0004-6361/202557450}{Astronomy \& Astrophysics (Mar. 2026)}},
  info={Equivalent preprint on \arxiv{2509.24740}},
]{SymBoltz.jl: A symbolic-numeric, approximation-free, and differentiable linear Einstein–Boltzmann solver}
\uioincludepdf{paper2.pdf}

\backmatter{}
\printbibliography{}

\end{document}
